\documentclass[svgnames]{article}    

\usepackage{graphicx}                 

% -- math                               
\usepackage{amssymb}
\usepackage{amsmath}
\usepackage{esint}
\usepackage{geometry}

% -- noindent
\setlength\parindent{0pt}

% -- images
%\graphicspath{{ }}                   

% -- tikz
\usepackage{pgfplots}
\pgfplotsset{compat=1.15}
\usepackage{comment}
\usetikzlibrary{arrows}
\usepackage[most]{tcolorbox}

% -- figures
\usepackage{float}
\usepackage{caption}
\usepackage{lipsum}

% -- start 
\title{Midterm 2 Review}
\author{Deval Deliwala}
%\date{}                              

\begin{document}
\maketitle
\tableofcontents

\vspace{20px}
\begin{center}
  
This document is not meant to cover all the math and derivations for key results in chapters 7, 8, and 9
of Griffiths (Perturbation Theory, Hydrogen Fine Structure, Variational Principle, WKB Approximation), 
but is just meant to outline the main results so I can recall them whenever.  

\end{center}

\newpage 

\section{Time Independent Perturbation Theory}

There are two types of time-independent perturbation theory. One for when the unperturbed energies 
of $\psi_n^{0}$ are non-degenerate (i.e. two different $\psi_n^{0}$, $\psi_m^{0}$ do not have the 
same energy. \\

The other is \textit{degenerate} perturbation theory, to account for when two unperturbed energy
eigenvalues are equivalent for different orthonormal wave functions. 

\subsection{Non-degenerate Perturbation Theory }

\subsubsection{First-Order Theory} 

Provided an unperturbed Schr\"odinger equation, 

\[
H^0 \psi_n^{0} = E_n^{0}\psi_n^0,       
\] \vspace{3px}

we obtain a \textit{complete set} of orthonormal eigenfunctions, $\psi_n^{0}$. 

\[
\langle \psi_n^{0}|\psi_m^{0} \rangle = \delta_nm.
\] \vspace{3px}


If we add a \textit{perturbation}, $H'$, modifying the Hamiltonian 

\[
H = H^{0} + H',
\] \vspace{3px}

We find the \textbf{first-order energy correction} to be equal to the expectation 
value of the perturbation hamiltonian with the \textit{unperturbed} eigenstate. That is, 

\[
    \boxed{
E_n^{1} = \langle \psi_n^{0} | H' | \psi_n^{0} \rangle .
}
\] \vspace{3px}

This is the fundamental result of first-order nondegenerate perturbation theory, and likely the most 
used equation of quantum mechanics. \\

The \textbf{first-order wave function correction} turns out to be equal to 

\[
    \boxed{
\psi_n^{1} = \sum_{m\neq n }^{} \frac{\langle \psi_m^{0} | H' |\psi_n^{0} \rangle}{E_n^{0} - E_m^{0}}\psi_m^{0}.
}\] \vspace{3px}

\subsubsection{Second-Order Theory} 

The second order correction to the eigenenergies in \textit{nondegenerate} perturbation theory 
is given by 

\[
\boxed{ 
\sum_{m\neq n}^{} \frac{|\langle \psi_m^{0}|H'|\psi_n^{0} \rangle |^2}{E_n^{0} - E_m^{0}}
}
\] \vspace{3px}

In practice the above equation is as far as it is useful to pursue perturbation theory for classical 
quantum mechanics. There is no need to calculate the second order correction to the wave function. 

\subsection{Degenerate Perturbation Theory}

If two unperturbed eigenenergies $E_n^{0}$ and $E_m^{0}$ are equivalent, you will notice in the denominator 
for the second-order energy corrections maps to 0. There is no reason to even trust the \textit{first-order} 
energy corrections. \\

Assume the unperturbed state takes the form 

\[
\psi^{0} = \alpha \psi_a^{0} + \beta \psi_b^{0}
\] \vspace{3px}


A perturbation usually breaks any energy degeneracy. That is, a degenerate energy splits into (say 2)  
nondegenerate energies after perturbation is added. \\

If we reverse the process, going from the 
splitted degeneracies back to the original degenerate state, the upper state reduces to one 
linear combination of $\psi_a^{0}$ and $\psi_b^{0}$ and the lower state reduces to an \textit{orthogonal}
linear combination -- but we do not know what these linear combinations \textit{are}. This is the reason 
we can't even use first-order nondegenerate theory. (I'll show a clever loophole around this later though). 

\subsection{First-Order Theory}

Let 

\[
    W_{ij} \equiv \big\langle \psi_i^{0} | H' | \psi_j^{0} \big\rangle
\] \vspace{3px}

where $(i,j) = \{a, b\}$. \\ 

The \textbf{first order energy correction} can be determined by solving the 
eigenvalue problem 

\[
\begin{pmatrix}
    W_{aa} & W_{ab} \\ 
    W_{ba} & W_{bb}
\end{pmatrix} 
\begin{pmatrix}
  \alpha \\ \beta
\end{pmatrix}
 = E^1 \begin{pmatrix}
     \alpha \\ \beta
 \end{pmatrix} 
\] \vspace{3px}

and the corresponding eigenvectors provide the coefficients $\alpha$ and $\beta$ that determine the 
``good states" (the reduced orthogonal linear combinations from reversing perturbations). Solving this problem 
yields the formula 

\[
\boxed{ 
    E_\pm^{1} = \frac{1}{2} \big[ W_{aa} + W_{bb} \pm \sqrt{(W_{aa} - W_{bb})^2 + 4|W_{ab}|^2}.
    }
\] \vspace{3px}

This is the fundamental result of degenerate perturbation theory. The two roots correspond to the two perturbed 
energies. 

\subsubsection{Clever Trick} 

Let $A$ be a hermitian operator that commutes with $H^0$ and $H'$ (a conserved quantity).
If $\psi_a^{0}$ and $\psi_b^{0}$ (the degenerate eigenfunctions of $H^{0}$ are \textit{also}
eigenfunctions of $A$, with \textit{unique} eigenvalues, 

\[
A\psi_a^{0} = \mu \psi_a^{0} \qquad A\psi_b^{0} = \nu \psi_b^{0} \qquad \mu \neq \nu,
\] \vspace{3px}

then $\psi_a^{0}$ and $\psi_b^{0}$ are the ``good states", and we can instead just use 
\textbf{nondegenerate} perturbation theory. \\

I won't show the proof here, but the moral is: \\ 

If you are faced with degenerate states, look around for some hermitian operator $A$ that commutes
with $H^{0}$ and $H'$. Pick your unperturbed states to be ones simultaneously eigenfunctions of 
$H_0$ and $A$ (with distinct eigenvalues). Then use \textit{ordinary} nondegerate perturbation theory. 

\subsection{Higher-Order Degeneracy} 

Previously I assumed the degeneracy was two-fold. In the case of an $n$-fold degeneracy, we look for eigenvalues 
of the $n\times n$ matrix 

\[
    W_{ij} = \Big\langle \psi_i^{0} | H' | \psi_j^{0}\Big\rangle
\] \vspace{3px}

And we similarly just solve for the eigenvalues of 

\[
\begin{pmatrix}
    W_{aa} & \hdots & W_{an} \\ 
    \vdots & \ddots & \vdots \\
    W_{na} & \hdots & W_{nn}
\end{pmatrix} \begin{pmatrix}
    \alpha \\ \vdots \\ \alpha_n
\end{pmatrix}  = E^1 \begin{pmatrix}
    \alpha \\ \vdots \\ \alpha_n 
\end{pmatrix} 
\] \vspace{3px}

which give the first-order energy corrections. And the good states are the corresponding eigenvectors 

\[
\psi^{0} = \alpha \psi_a^{0} + \cdots  + \alpha_n \psi_n^{0}. 
\] \vspace{3px}

\section{The Fine Structure of Hydrogen} 

When calculating the $n$th wave function of hydrogen, we used the \textit{Bohr} Hamiltonian, 

\[
H_\text{bohr} = -\frac{\hbar^2}{2m}\nabla^2 - \frac{e^2}{4\pi\epsilon_0} \frac{1}{r}
\] \vspace{3px}

We corrected for the motion of the nucleus by replacing $m$ by the reduced mass $\mu$. More significant is 
the \textbf{fine structure} which arises due to a \textit{relativistic correction} and \textit{spin-orbit coupling}.
Fine structure gives a tiny perturbation -- smaller than a factor of $\alpha^2$, where $\alpha$ is given by 

\[
\alpha \equiv \frac{e^2}{4\pi\epsilon_0\hbar c} \approx \frac{1}{137.036},
\] \vspace{3px}

the famou \textbf{fine-structure constant}. 

\subsection{Relativistic Correction} 

Relativistic momentum is given by 

\[
    p = \frac{mv}{\sqrt{1 - (v/c)^2}}
\] \vspace{3px}

And relativistic kinetic energy is given by 

\[
    T = \frac{mc^2}{\sqrt{1 - (v/c)^2}}
\] \vspace{3px}

Rearranging gives a simpler formula for $T$, 

\[
    T = \sqrt{p^2c^2 + m^2 c^4} - mc^2
\] \vspace{3px}

Expanding into a taylor series, using the fact that 

\[
    \sqrt{1+x} \approx \left(1 + \frac{1}{2}x - \frac{1}{8}x^2 + \cdots\right)
\] \vspace{3px}

We get 

\begin{align*}
    T &= mc^2\left[ 1 + \frac{1}{2}\left( \frac{p}{mc}  \right) ^2 - \frac{1}{8} \left( \frac{p}{mc} \right) ^4 \cdots - 1 \right]\\
      &= \frac{p^2}{2m} - \frac{p^4}{8m^3c^2} + \cdots.
\end{align*} \vspace{3px}


The lowest-order relativistic correction to the Hamiltonian is therefore 

\[
H_r' = -\frac{p^4}{8m^3 c^2}
\] \vspace{3px}

Where we can use first-order perturbation theory to calculate the first-order energy correction to $E_n$, 

\[
E_r^{1} = \langle H_r' \rangle = -\frac{1}{8m^3c^2}\langle \psi \rangle p^4 \psi \rangle = -\frac{1}{8m^3c^2}\langle p^2 \psi | p^2 \psi \rangle 
\] \vspace{3px}

But the Schr\"odinger equation says that (for unperturbed states) 

\[
p^2\psi = 2m(E-V)\psi
\] \vspace{3px}

Using this fact and simplifying we can calculate the general relativistic correction for any potential $V$ to be 

\[
    \boxed{
    E_r^{1} = -\frac{1}{2mc^2} \Big\langle (E-V)^2 \Big\rangle = -\frac{1}{2mc^2} \left[ E^2 - 2E\langle V\rangle + \langle V^2 \rangle \right].
}\] \vspace{3px}

which is a very small correction (on the order of $1/c^2$). The bigger correction is generated via  \textit{spin-orbit coupling}. 

\subsection{Spin-Orbit Coupling} 

Imagine the electron is in orbit around the nucleus. From the electron's POV, it's actually the proton circling around \textit{it}. 
The orbiting positive charge generates a magnetic field $B$ in the electron frame, which thus exerts a torque 
on the spinning electron, tending to align its magnetic moment $\mu$ along the direction of the field. 
The Hamiltonian is given by 

\[
H = -\vec{v} \cdot \vec{B}
\] \vspace{3px}

where $|B| = \frac{\mu_0 I}{2r}$. And $I$ is the effective current $I = e/T$, where $T$ is the period of the orbit. 
Combining this, with the orbital angular momentum of the electron (in the rest frame of the nucleus), 
we can calculate 

\[
\vec{B} = \frac{1}{4\pi \epsilon_0} \frac{e}{mc^2r^3}\vec{L}
\] \vspace{3px}

where we replaced $c = \frac{1}{\sqrt{\epsilon_0 \mu_0}}$. If we include a spin and relativistic correction, we 
get a modified Hamiltonian (accounting for Thomas precession), 

\[
    H_{so}' = \left( \frac{e^2}{8\pi\epsilon_0} \right) \frac{1}{m^2c^2r^3}\vec{S} \cdot \vec{L} ,
\] \vspace{3px}

The \textbf{spin-orbit} interaction.  Doing a bunch of calculations we get the energy correction to be 

\[
E_{so}^{1} = \frac{E_n^2}{mc^2} \left\{ \frac{n \left[ j(j+1)-\ell(\ell+1) - 3/4 \right] }
{\ell(\ell+1/2)(\ell+1)} \right\}
\] \vspace{3px}


Adding this to the relativistic correction, we get the complete fine-structure formula, 

\[
\boxed{ 
    E_{fs}^{1} = \frac{(E_n)^2}{2mc^2} \left( 3 - \frac{4n}{j + 1/2} \right) .  
}
\] \vspace{3px}

\subsubsection{Zeeman Effect} 

When an atom is placed in a uniform external magnetic field $\vec{B}_\text{ext} $, the energy levels get shifted. \\

For a single electron, 

\[
H'_Z = -(\mu_l + \mu_s) \cdot \vec{B}_\text{ext} 
\] \vspace{3px}

where  $\mu_s = -\frac{e}{m}\vec{S}$ and $\mu_l = -\frac{e}{2m}\vec{L}$. Hence, 

\[
H'_Z = \frac{e}{2m}(\vec{L}+2\vec{S}) \cdot \vec{B}_\text{ext} . 
\] \vspace{3px}

\subsubsection{Hyperfine Splitting in Hydrogen} 

The proton also is a magnetic dipole, though its dipole moment is much smaller than the electron's 
because of the mass in the denominator 

\[
\mu_p = \frac{g_p e}{2m_p}\vec{S}_p, \quad \mu_e = -\frac{e}{m_e}\vec{S}_e
\] \vspace{3px}

The $g$-factor comes from the fact the proton is made up of three quarks. There's stuff that comes as 
a result of all this, specifically  \textit{spin-spin}  coupling between the proton and electron but 
we ultimately get a \textit{hyperfine energy correction} which is given by 

\[
    E_\text{hf}^1 = \frac{4g_p \hbar^{4}}{3m_p m_e^2 c^2 a_4} \begin{cases} +1/4 \text{ (if triplet)},  \\ -3/4 \text{ (if singlet)}  \end{cases}
\] \vspace{3px}












\end{document}
